% ══════════════════════════════════════════════════════
%  contenu/exos-problemes.tex  —  Exercices et problèmes
% ══════════════════════════════════════════════════════
\clearpage
\PartieHeader{Exercices et probl\`emes}

\begin{multicols*}{2}

%% ─── Thème 1 : ED du premier ordre y'=ay ──────────
\ThemeHeader{ED du premier ordre $y'=ay$}

\begin{Exercice}
R\'esoudre les \'equations diff\'erentielles suivantes :
\begin{enumerate}[label=\alph*)]
  \item $y'+3y=0$
  \item $y'-\sqrt{2}\,y=0$
  \item $2y'+3y=0$
  \item $y'+ey+e^2=ey(e^2-1)$
\end{enumerate}
\end{Exercice}

\begin{Exercice}
D\'eterminer la solution particuli\`ere de chaque \'equation diff\'erentielle :
\begin{enumerate}[label=\alph*)]
  \item $y'+3y=0$ v\'erifiant $y(0)=-5$
  \item $y'-5y=0$ v\'erifiant $y(\ln 2)=5$
  \item $y'-\sqrt{2}\,y=0$ v\'erifiant $y(-1)=0$
  \item $y'+ey=0$ v\'erifiant $y(e)=0$
\end{enumerate}
\end{Exercice}

\begin{Exercice}
V\'erifier que les solutions de l'\'equation diff\'erentielle
$(E):y'-y\ln 2=0$ sont les fonctions de la forme
$x\mapsto\alpha\cdot 2^x$, $\alpha$ un r\'eel.
\end{Exercice}

%% ─── Thème 2 : ED du premier ordre y'=ay+b ────────
\ThemeHeader{ED du premier ordre $y'=ay+b$}

\begin{Exercice}
R\'esoudre les \'equations diff\'erentielles suivantes et d\'eterminer la solution v\'erifiant la condition initiale donn\'ee :
\begin{enumerate}[label=\alph*)]
  \item $y'+3y=1$ et $y(0)=0$
  \item $y'=5y$ et $y(1)=-2$
  \item $y'+y+2=0$ et $y(0)=0$
  \item $y'+3y=1$ et $y(0)=0$
  \item $y'+ey+e^2=et$ et $y(e)=0$
  \item $y'+ey+e^2=ey\,(e^2-1)$
\end{enumerate}
\end{Exercice}

\begin{Exercice}
\begin{enumerate}
  \item D\'eterminer la solution g\'en\'erale de $y'+7y+11=0$.
  \item D\'eterminer la solution de $y'=\pi y+\sqrt{2}$ v\'erifiant $y(0)=-1$.
\end{enumerate}
\end{Exercice}

\begin{Exercice}
\begin{enumerate}
  \item R\'esoudre $(E):y'-2y=e^{2x}$.

  \item Soit $f$ d\'erivable sur $\R$ telle que $f'(x)-2f(x)=e^{2x}$.
    \begin{enumerate}[label=\alph*)]
      \item V\'erifier que $g(x)=xe^{2x}$ est solution de $(E)$.
      \item R\'esoudre $(E')$ : $y'-2y=0$.
      \item Montrer que $f$ est solution de $(E)$ ssi $(f-g)$ est solution de $(E')$.
      \item D\'eduire toutes les solutions de $(E)$.
      \item D\'eterminer la solution $h$ de $(E)$ v\'erifiant $h(0)=1$.
    \end{enumerate}
\end{enumerate}
\end{Exercice}

\begin{Exercice}
Soit l'\'equation diff\'erentielle $(E):y'-4y=-8x^2$.
\begin{enumerate}
  \item Chercher une solution particuli\`ere $P$ polyn\^ome du second degr\'e.
  \item Poser $y_1=y-P$. Montrer que $y$ est solution de $(E)$ ssi $y_1$ v\'erifie $y_1'-4y_1=0$.
  \item En d\'eduire toutes les solutions de $(E)$.
\end{enumerate}
\end{Exercice}

\begin{Exercice}
On consid\`ere les deux \'equations :
$(E_0):y'-2y=0$ \quad et \quad $(E):y'-2y=2(e^{2x}-1)$.
\begin{enumerate}
  \item R\'esoudre $(E_0)$.
  \item V\'erifier que $f(x)=2xe^{2x}+1$ est solution de $(E)$.
  \item Montrer que $(g-f)$ est solution de $(E_0)$ ssi $g$ est solution de $(E)$.
  \item En d\'eduire toutes les solutions de $(E)$.
\end{enumerate}
\end{Exercice}

%% ─── Thème 3 : ED du second ordre ─────────────────
\ThemeHeader{ED du second ordre $y''+ay'+by=0$}

\begin{Exercice}
R\'esoudre les \'equations diff\'erentielles suivantes :
\begin{enumerate}[label=\alph*)]
  \item $y''+4y=0$
  \item $y''-y'=0$
  \item $y''-2y=0$
  \item $y''+y'+y=0$
  \item $y''+4y'-5y=0$
\end{enumerate}
\end{Exercice}

\begin{Exercice}
R\'esoudre les \'equations diff\'erentielles suivantes et d\'eterminer la solution v\'erifiant les conditions initiales :
\begin{enumerate}[label=\alph*)]
  \item $y''-5y'+6y=0$\,;\, $y(0)=-1$, $y'(0)=1$
  \item $y''+y'+y=0$\,;\, $y(0)=0$, $y'(0)=0$
  \item $3y''+2y'+y=0$\,;\, $y(\pi\sqrt{2})=2$, $y'(\pi\sqrt{2})=0$
  \item $y''+y'+y=0$\,;\, $y(0)=0$, $y'(0)=0$
\end{enumerate}
\end{Exercice}

\begin{Exercice}
\begin{enumerate}
  \item R\'esoudre $(E):y''-2y'+5y=0$.
  \item D\'eterminer la solution $f$ de $(E)$ v\'erifiant $f(0)=1$ et $f'(0)=1$.
  \item D\'eterminer une primitive de $f$.
\end{enumerate}
\end{Exercice}

\begin{Exercice}
R\'esoudre et discuter selon les valeurs de $m$ :
$(E_m):y''-2my'+my=0$.
\end{Exercice}

\begin{Exercice}
On consid\`ere les deux \'equations diff\'erentielles :
$(1):y''-2y'+3y=0$ \quad et \quad $(2):z''+2z'=0$.
On pose $y=ze^x$ pour tout $x\in\R$.
\begin{enumerate}
  \item Montrer que $y$ est solution de $(1)$ ssi $z$ est solution de $(2)$.
  \item D\'eterminer les solutions de $(2)$.
  \item D\'eterminer $f$, solution de $(1)$, v\'erifiant $f(0)=1$ et $f'(0)=-2$.
\end{enumerate}
\end{Exercice}

%% ─── Thème 4 : Applications physiques ─────────────
\ThemeHeader{Applications physiques}

\begin{Exercice}
\textbf{Circuit RC.}
On consid\`ere le circuit \'electrique d\'ecrit par :
$R\dfrac{dq}{dt}+\dfrac{q}{C}=U$ avec $q(0)=0$.
\begin{enumerate}
  \item D\'eterminer toutes les solutions de l'\'equation.
  \item Montrer que $q(t)=CU\!\left(1-e^{-t/(RC)}\right)$.
  \item On observe la tension $U_r=R\cdot i$ aux bornes du r\'esistor.
        Sachant que $i=\dfrac{dq}{dt}$, donner l'expression de $U_r$ en fonction de $t$.
\end{enumerate}
\end{Exercice}

\begin{Exercice}
\textbf{Croissance bact\'erienne.}
La vitesse d'accroissement des bact\'eries est proportionnelle
au nombre de bact\'eries pr\'esentes.

On note $N(t)$ le nombre de bact\'eries (en millions). La vitesse v\'erifie :
$N'(t)=kN(t)$, avec $N(0)=N_0$.
\begin{enumerate}
  \item Donner l'expression de $N(t)$.
  \item En combien de temps le nombre de bact\'eries sera-t-il doubl\'e ?
\end{enumerate}
\end{Exercice}

\begin{Exercice}
\textbf{Dynamique de population (mod\`ele logistique).}
On note $u(t)$ le nombre de rongeurs dans une r\'egion. On suppose que $u$ v\'erifie :
\[
  (E_2):\begin{cases}u'(t)=\dfrac{u(t)}{4}-\dfrac{[u(t)]^2}{12}\\u(0)=1\end{cases}
\]
\begin{enumerate}[label=\alph*)]
  \item On suppose $u(t)>0$ et on pose $h(t)=\dfrac{1}{u(t)}$.
        Montrer que $u$ v\'erifie $(E_2)$ ssi $h$ v\'erifie :
        $(E_3):\begin{cases}h'=-\dfrac{h}{4}+\dfrac{1}{12}\\h(0)=1\end{cases}$
  \item R\'esoudre $(E_3)$ et en d\'eduire $u(t)$.
  \item Comment se comporte la population lorsque $t\to+\infty$ ?
\end{enumerate}
\end{Exercice}

%% ─── Thème 5 : Problème avancé ─────────────────────
\ThemeHeader{Probl\`eme avanc\'e}

\begin{Exercice}
\textbf{\'Equation d'ordre 3.}

Soit l'\'equation diff\'erentielle $(E):y'''-6y''+12y'-8y=0$.
\begin{enumerate}
  \item V\'erifier que $h:x\mapsto e^{2x}$ est solution de $(E)$.
  \item Soit $f$ une fonction trois fois d\'erivable sur $\R$ et $g:x\mapsto f(x)e^{-2x}$.
        Montrer que $f$ est solution de $(E)$ si et seulement si $g'''$ est la fonction nulle.
  \item En d\'eduire toutes les solutions de $(E)$.
\end{enumerate}
\end{Exercice}

\begin{Exercice}
\textbf{\'Equation avec second membre.}

Soit $(E):y''+y'=\dfrac{e^x}{x^2}$ sur $]0;+\infty[$.

On donne $(E_0):y'+y=0$.
\begin{enumerate}
  \item R\'esoudre $(E_0)$.
  \item Montrer que $f(x)=\dfrac{e^x}{x}$ est solution de $(E)$.
  \item Montrer que $(g-f)$ est solution de $(E_0)$ sur $]0;+\infty[$ ssi $g$ est solution de $(E)$.
  \item D\'eduire toutes les solutions de $(E)$ sur $]0;+\infty[$.
\end{enumerate}
\end{Exercice}

\begin{Exercice}
\textbf{\'Equation $4y''+9y=0$. Probl\`eme g\'eom\'etrique.}
\begin{enumerate}
  \item R\'esoudre l'\'equation diff\'erentielle $(E):4y''+9y=0$.
  \item On d\'esigne par $f$ la solution particuli\`ere de $(E)$ dont la courbe repr\'esentative
        admet une tangente parall\`ele \`a l'axe des abscisses au point $A\!\left(\dfrac{\pi}{6};2\right)$.
    \begin{enumerate}[label=\alph*)]
      \item D\'eterminer une expression de $f(x)$.
      \item Montrer que pour tout r\'eel $x$ :
            $f(x)=2\cos\!\left(\dfrac{3x}{2}-\dfrac{\pi}{4}\right)$.
    \end{enumerate}
\end{enumerate}
\end{Exercice}

\end{multicols*}
